
LLM-as-judge is increasingly deployed to evaluate code correctness without test execution, yet practitioners lack guidance on how model scale affects reliability. This paper presents a scale-controlled analysis of error patterns in LLM code judges, comparing 7B, 70B, and proprietary models under fixed evaluation settings. The central finding is that scale predicts error \textit{type}, not merely error rate: smaller models systematically over-accept incorrect code (FPR=74.8\%), while larger models systematically under-accept correct code (FNR=29.3\%). This FPR-FNR tradeoff is statistically significant ($\chi^{2}=45.78$, p=3.$27\times 10^{-8})$. Two predictions based on the complementary-errors hypothesis were falsified: (1) ensemble voting across scales degrades accuracy by 8.05\% compared to the best single judge (McNemar p=1.$45\times 10^{-6})$, and (2) unanimous agreement correlates with lower accuracy (35.8\%) than disagreement (39.7\%). Judge-execution agreement increases with scale (7B: 58.5\%, 70B: 70.7\%, proprietary: 71.3\%) with diminishing returns (20:1 ratio at the 70B transition, Kruskal-Wallis p=0.021). These results suggest that scale selection functions as error-type selection, and practitioners should choose judges based on error-cost asymmetry rather than ensemble combination.

